\section{Discussion and Limitations}
\label{sec:discussion}

\subsection{What the Hellman experiment says}

The Hellman sweep began from an apparent discrepancy: at the configuration
$mt=N$, an independent-occupancy heuristic suggests about $63.2\%$ coverage,
while the original experiment measured only about $16.8\%$. Once the correct
fixed-map baseline is used, the discrepancy disappears. The Ma--Hong
recurrence predicts \StressMH{}, and both the DES-derived mapping and matched
random maps are close to that value.

This is a useful negative result. It prevents an implementation artifact or an
inappropriate iid approximation from being promoted into a cryptanalytic
claim. Within the state dimensions tested here, we find no stable Hellman
coverage anomaly that distinguishes the DES-derived mapping from the random
controls.

\subsection{What the distinguished-point experiment adds}

Distinguished points expose a different finite-functional-graph effect. The
fresh-sample model predicts an exponentially decreasing truncation probability
$(1-2^{-d})^L$. For a fixed mapping, however, a trajectory eventually enters a
cycle. If the relevant eventual cycle does not contain a distinguished state,
increasing $L$ beyond the transient and cycle lengths cannot make the chain
terminate.

The exact reverse-reachability computation quantifies this phenomenon over all
$2^{16}$ starts. In particular, the large-$L$ truncation plateau observed for
$d=8$ is not a timing artifact or a failure of the DP implementation: it is a
property of the complete functional graph. The classical DP literature already
treats bounded chains and cycle removal as practical concerns
~\cite{StandaertEtAl2002,AvoineJunodOechslin2008}; our contribution here is
the exact executable finite-state accounting for the concrete map and matched
random controls.

\subsection{Why reduced DES is useful here}

DES is obsolete and the experiments do not attempt to recover a full DES key.
The reduced model is useful precisely because exact enumeration is possible.
For $b\leq18$ we can construct the entire iteration map, measure exact Hellman
coverage, and compare with complete random mappings. At $b=16$ we can also
compute exact reverse reachability for every distinguished-point predicate.

The resulting object is best viewed as a finite cryptanalytic laboratory: the
primitive supplies a concrete deterministic mapping, while the reduced
dimension makes functional-graph quantities observable without statistical
estimation.

\subsection{Limitations}

Several limitations delimit the claims. First, the DES-derived mapping depends
on one fixed plaintext, one key embedding, and one reduction. Different
plaintexts or reductions induce different functional graphs. Second, the
Hellman study stops at $b=18$ and the exact distinguished-point graph analysis
is performed at $b=16$; no asymptotic statement about larger DES subspaces
follows automatically. Third, the random controls are deterministic
SplitMix64-generated mappings used for reproducibility. They are experimental
stand-ins for uniformly sampled functions, not cryptographic random oracles.

Fourth, the bootstrap intervals in Section~\ref{sec:results} are descriptive
summaries of the chosen 30 paired trials. We do not present them as a formal
proof of equivalence between the DES-derived map and the random-function
distribution. Fifth, the Java/Rust result establishes semantic equality on a
shared workload, not equality of all possible implementations or providers.

Finally, no Java-versus-Rust speedup claim appears in this version. Although
both languages contain timing harnesses, the scientific contribution of this
paper is exact finite-state behavior and reproducibility. A performance study
requires a separately controlled benchmarking design.

\subsection{Implications for TMTO experimentation}

The experiments suggest three practical rules for reproducible TMTO studies:
use the coverage model appropriate to iterated fixed functions rather than iid
occupancy; separate exact semantic validation from performance measurement;
and, for distinguished points, distinguish short-chain geometric behavior from
the eventual reachability structure of the functional graph.
